\documentclass[11pt]{article}

\usepackage[T1]{fontenc}   % proper glyph encoding so fi/fl/quote ligatures render
\usepackage{lmodern}
\usepackage[margin=1in]{geometry}
\usepackage{amsmath,amssymb}
\usepackage{enumitem}
\usepackage{booktabs}
\usepackage[colorlinks=true,linkcolor=blue,citecolor=blue,urlcolor=blue]{hyperref}

\newcommand{\GF}{\mathbb{F}_2}
\newcommand{\rank}{\operatorname{rank}}
\newcommand{\wt}{\operatorname{wt}}
\newcommand{\eff}{kd^2/n}

\title{\textbf{The qLDPC Challenge}\\[4pt]
  \large A public, automatically verified leaderboard for quantum LDPC codes}
\author{Unitary Foundation\\
  \texttt{github.com/unitaryfoundation/qldpc-challenge}}
\date{Extended abstract, lightning talk, Quantum Software Workshop 2026
  (IEEE Quantum Week, Toronto)}

\begin{document}
\maketitle

\begin{abstract}
Quantum low-density parity-check (qLDPC) codes are the leading route to
low-overhead fault tolerance, but the field has no shared, machine-checked record
of which codes are best under which constraints. Parameters are scattered across
papers and tables, and distance claims in particular are hard to reproduce and
easy to overstate, since computing distance is NP-hard and is often estimated
with heuristic decoders that can be wrong. The qLDPC Challenge is a public
leaderboard that treats this as a software-infrastructure problem. A participant
submits one JSON file describing a CSS code; continuous integration recomputes
the code's properties from its parity-check matrices, checks a self-certifying
distance witness, and merges the entry only if it holds up. Ranking is a
per-track Pareto frontier rather than a single gameable number, track membership
is computed rather than self-declared, and the whole pipeline (validator,
distance refutation gate, exact certifier, scorer, static-site generator) lives
in the repository and runs on every pull request. We describe the trust model,
the verification pipeline, and the current board, and argue the design is a
reusable pattern for reproducible, self-certifying benchmarks in quantum error
correction.
\end{abstract}

\section{Motivation}

Recent qLDPC constructions, from the bivariate-bicycle ``gross'' code
$[[144,12,12]]$ of Bravyi et al.\ to planar open-boundary families, have made
order-of-magnitude reductions in the qubit overhead of fault tolerance plausible.
Progress is now fast and distributed across many groups, which makes a shared
question pressing: for a given hardware constraint (check weight, geometric
locality, code size), what is the best known code, and can we trust the claim?

Two facts make this hard to answer in the literature alone. First, results live
in heterogeneous tables across papers, with no common format or single place to
compare like with like. Second, and more seriously, the headline parameter,
distance $d$, is the one that is hard to verify. Finding a low-weight logical
operator gives only an upper bound on $d$; proving $d \ge D$ is the NP-hard
direction. Reported distances are frequently decoder-based upper bounds, and
recent work documents belief-propagation decoders overestimating qLDPC distances
by large factors. A leaderboard that simply trusts submitted numbers would
accumulate quiet errors.

The qLDPC Challenge addresses both problems with quantum-software tooling rather
than a static document. Every entry is a small machine-readable file that is
automatically re-verified, distance claims must carry a checkable witness, and a
distance refutation step actively searches for counterexamples. The result is a
reproducible, community-extensible, self-certifying record.

\section{Design}

Three principles shape the system.

\paragraph{(i) Everything cheap is trustless and mandatory.} The verifier
recomputes $n$, the logical dimension
$k = n - \rank_{\GF}(H_X) - \rank_{\GF}(H_Z)$, CSS commutation
$H_X H_Z^{\top} = 0$ over $\GF$, the maximum check weight, and the geometric
locality class, all from the submitted parity checks. A wrong claimed $k$ or a
non-commuting ``stabilizer'' is rejected in milliseconds. Nothing in this tier is
taken on trust.

\paragraph{(ii) Distance is layered by how much is proven.} Because exact
distance does not scale, the board separates three confidence tiers. An
\emph{upper bound} ($d \le$) is backed by an explicit logical-operator witness
that the verifier confirms is a genuine nontrivial logical of the claimed weight,
so the bound is checkable with no trust required. A \emph{corroborated} bound
($d \le^{\!*}$) is an upper bound that an independent search has tried and failed
to beat. An \emph{exact} value ($d =$) is one a server-side integer program has
proven minimal. Nothing is ever silently upgraded.

\paragraph{(iii) A code is not a single number.} The primary tracks are a grid
of two nested, \emph{computed} axes: a locality class
(\texttt{local-2d-single} $\subset$ \texttt{local-2d-bilayer} $\subset$
\texttt{unrestricted}, derived from the layout) and a check-weight class
(\texttt{weight-4} $\subset$ \texttt{weight-6} $\subset$ \texttt{weight-8}
$\subset$ any). Membership is derived from the parity checks and the layout, so a
track cannot be claimed by relabeling. Within each cell the ranking is a Pareto
frontier over $(n,k,d,w)$; there is no single winner. The construction family
(bivariate bicycle, generalized bicycle, two-block group algebra, tile,
topological, ...) cannot be recovered from the matrices, so it is a filterable
tag, never a ranking.

\section{The verification pipeline}

The pipeline is the reusable artifact. On every pull request, GitHub Actions
runs, in order:

\begin{enumerate}[label=\arabic*.,leftmargin=*,nosep]
  \item \textbf{Structural and trustless checks.} Schema validity, qubit indices
  in range, no repeated index within a check, CSS commutation, recomputed $k$
  matching the claim, $k \ge 1$, a controlled family and novelty vocabulary, and
  the witness check: each distance witness must have weight equal to the claimed
  value, lie in the kernel of the opposite checks, and be a nontrivial logical
  (not a stabilizer product). This certifies the per-side upper bound.
  \item \textbf{Computed locality class.} If a layout is present, the verifier
  derives the locality class from the coordinates, enforcing no cramming (at most
  \texttt{layers} qubits per site, distinct sites at least unit-spaced) and a
  bounded interaction radius, so a small radius cannot be faked by collapsing
  qubits.
  \item \textbf{Distance refutation gate.} Two independent searches, a
  random-information-set search and a syndrome-decoder search, look for a logical
  lighter than the claimed distance. Any hit is a sound refutation (it exhibits a
  checkable counterexample). The per-run seed is random by design so an
  over-claim cannot be tuned to evade one fixed search, and the gate fails closed.
  A weekly whole-board sweep re-runs the refutation with fresh seeds to
  accumulate coverage over time.
  \item \textbf{Identity and duplicate checks.} The pull-request author must be a
  listed author of the codes they add. Every code is fingerprinted two ways: the
  reduced-row-echelon form of $(H_X, H_Z)$ pins the exact stabilizer group (an
  identical fingerprint is a hard error), and a permutation-invariant
  Weisfeiler--Leman signature on the Tanner graph flags possible equivalents for
  review.
\end{enumerate}

A separate, maintainer-run exact tier upgrades a claim to $d =$. For each logical
generator $t_L$ it asks, via a bounded integer program (scipy/HiGHS, no
commercial solver), whether a $v$ exists with $H v \equiv 0$,
$\langle v, t_L\rangle \equiv 1$, and $\wt(v) \le \texttt{value} - 1$. If every
generator on both sides is infeasible, the distance is exact; if a lighter
logical is found, the claim is refuted; if the solver times out, the side
honestly remains an upper bound. A certificate written to \texttt{certs/} is what
promotes the board label.

\section{Metric and the state of the board}

The figure of merit the 2D-locality literature uses, $\eff$, is the
Bravyi--Poulin--Terhal saturation ratio. For a 2D-local or bounded-weight code it
is bounded by a constant that depends only on the interaction range, check
weight, and on-site dimension, not on $n$, so maximizing it within a fixed
locality model is well posed. For asymptotically good codes, where $k$ and $d$
both grow with $n$, it grows like $n^2$ without bound. The board therefore treats
$\eff$ as a per-cell, sortable column and compares like with like, never as a
global record. Table~\ref{tab:ref} lists reference points.

\begin{table}[h]
\centering
\begin{tabular}{lcccc}
\toprule
Code & $[[n,k,d]]$ & topology & $w$ & $\eff$\\
\midrule
Rotated surface (Kitaev)      & $[[d^2,1,d]]$   & planar   & 4 & $1$\\
Gross code (Bravyi et al.)    & $[[144,12,12]]$ & periodic & 6 & $12$\\
Tile code (Eberhardt et al.)  & $[[512,18,19]]$ & planar   & 8 & $12.7$\\
Panteleev--Kalachev (GB)      & $[[126,28,8]]$  & periodic & 10 & $14.2$\\
\bottomrule
\end{tabular}
\caption{Reference points; $\eff$ values for 2D-local weight-8 codes are
exact, integer-programming-certified.}
\label{tab:ref}
\end{table}

At the time of writing the board holds 40 verified codes, 21 submitted through
the challenge and 19 literature baselines, with 27 certified exact, across the
bivariate-bicycle, generalized-bicycle, two-block group-algebra, tile, and
topological families. It is seeded against the published state of the art (the
planar codes of Liang, Eberhardt, and Chen; the bivariate-bicycle codes of Bravyi
et al.; the generalized-bicycle codes of Panteleev--Kalachev and Wang--Pryadko;
and the Kitaev and Steane baselines), so a submission is always measured against
known codes. The seed set is selected, not an exhaustive snapshot, and the board
states this: an on-board leader may still be matched by a published code that has
not been seeded. The badges and figures are regenerated from the live board on
every build, so they track the current numbers rather than a hand-typed count.

\section{Relevance to quantum software, and the talk}

The challenge is a small case study in software infrastructure for quantum error
correction. It encodes a trust model directly in CI: cheap properties are
recomputed and mandatory, distance is self-certifying at the cheap tier and
separately earned at the exact tier, and an adversarial refutation step guards
the one quantity that is expensive to prove. The submission unit is a verifiable
file, the referee is a script anyone can run locally, and the contribution
workflow is a pull request, which lowers the barrier to comparing codes while
keeping over-claims out of the record. The verifier and the refutation gate are
reusable beyond this board.

The lightning talk will demonstrate the end-to-end flow (submit, CI verifies,
board updates), the computed layered scheme and its trust model, and how to
contribute, including via LLM-guided search over code-generating programs. We
will also raise the questions the board is built to surface: whether planar
weight-6 codes can close the gap to periodic records, how far the refutation gate
and exact certification scale, and which families have unmined finite-size
headroom.

\subsection*{Key references}
\begin{itemize}[nosep]
  \item Bravyi, Cross, Gambetta, Maslov, Rall, Yoder, Nature \textbf{627}, 778
  (2024), arXiv:2308.07915 (bivariate-bicycle / gross code).
  \item Liang, Eberhardt, Chen, arXiv:2504.08887 (planar open-boundary qLDPC).
  \item Eberhardt et al., tile codes, arXiv:2504.09171.
  \item Panteleev, Kalachev, Quantum \textbf{5}, 585 (2021), arXiv:1904.02703;
  Wang, Pryadko, arXiv:2203.17216 (generalized bicycle).
  \item Bravyi, Poulin, Terhal, Phys.\ Rev.\ Lett.\ \textbf{104}, 050503 (2010),
  arXiv:0909.5200 (the $\eff$ bound).
\end{itemize}

\end{document}
